\textbf{Phần 1. Trắc nghiệm lựa chọn}
\setcounterpageref{ex}{0}

\begin{ex} %[2D11V55]
$1\text{ m}$ bằng bao nhiêu đề-xi-mét?
\choice
{$100\text{ dm}$}
{\True $10\text{ dm}$}
{$1\text{ dm}$}
{$1 000\text{ dm}$}
\loigiai{
$1\text{ m} = 10\text{ dm}$.
\textbf{Đáp án của câu hỏi là B}
}
\end{ex}

\begin{ex} %[2D11V55]
$1\text{ km}$ bằng bao nhiêu mét?
\choice
{$10\text{ m}$}
{$100\text{ m}$}
{$10 000\text{ m}$}
{\True $1 000\text{ m}$}
\loigiai{
$1\text{ km} = 1 000\text{ m}$.
\textbf{Đáp án của câu hỏi là D}
}
\end{ex}

\begin{ex} %[2D11V55]
Độ dài chiếc bút chì thông thường khoảng bao nhiêu?
\choice
{\True $15\text{ cm}$}
{$15\text{ m}$}
{$15\text{ km}$}
{$15\text{ dm}$}
\loigiai{
Bút chì dài khoảng 15 cm.
\textbf{Đáp án của câu hỏi là A}
}
\end{ex}

\begin{ex} %[2D11V55]
Một sợi dây dài $2\text{ m}$. Người ta cắt đi $6\text{ dm}$. Chiều dài sợi dây còn lại là
\choice
{$16\text{ dm}$}
{$4\text{ dm}$}
{\True $14\text{ dm}$}
{$14\text{ m}$}
\loigiai{
Đổi $2\text{ m} = 20\text{ dm}$. Còn lại: $20 - 6 = 14$ dm.
\textbf{Đáp án của câu hỏi là C}
}
\end{ex}

\begin{ex} %[2D11V55]
Một quãng đường gồm hai chặng: chặng một dài $400\text{ m}$, chặng hai dài $600\text{ m}$. Cả hai chặng đường dài bao nhiêu ki-lô-mét?
\choice
{$10\text{ km}$}
{\True $1\text{ km}$}
{$2\text{ km}$}
{$100\text{ km}$}
\loigiai{
$400 + 600 = 1 000\text{ m} = 1\text{ km}$.
\textbf{Đáp án của câu hỏi là B}
}
\end{ex}

\begin{ex} %[2D11V55]
Điền số thích hợp: $3\text{ m } 5\text{ dm} = \dots\text{ dm}$.
\choice
{$305\text{ dm}$}
{$8\text{ dm}$}
{$53\text{ dm}$}
{\True $35\text{ dm}$}
\loigiai{
$3\text{ m} = 30\text{ dm} \Rightarrow 3\text{ m } 5\text{ dm} = 35$ dm.
\textbf{Đáp án của câu hỏi là D}
}
\end{ex}

\begin{ex} %[2D11V55]
Con đường từ nhà Mai đến trường dài $500\text{ m}$. Mai đi từ nhà đến trường rồi lại đi từ trường về nhà. Mai đã đi tất cả bao nhiêu mét?
\choice
{\True $1 000\text{ m}$}
{$500\text{ m}$}
{$800\text{ m}$}
{$100\text{ m}$}
\loigiai{
$500 + 500 = 1 000$ (m).
\textbf{Đáp án của câu hỏi là A}
}
\end{ex}

\begin{ex} %[2D11V55]
Điền dấu thích hợp: $4\text{ m} \dots 38\text{ dm}$.
\choice
{$<$}
{$=$}
{\True $>$}
{Không so sánh được}
\loigiai{
$4\text{ m} = 40\text{ dm} > 38\text{ dm}$.
\textbf{Đáp án của câu hỏi là C}
}
\end{ex}

\begin{ex} %[2D11V55]
Bác thợ có thanh gỗ dài $1\text{ m}$. Bác cưa thành các đoạn ngắn dài $2\text{ dm}$. Bác cưa được bao nhiêu đoạn?
\choice
{$10$ đoạn}
{\True $5$ đoạn}
{$2$ đoạn}
{$8$ đoạn}
\loigiai{
$1\text{ m} = 10\text{ dm}$. Số đoạn: $10 : 2 = 5$ đoạn.
\textbf{Đáp án của câu hỏi là B}
}
\end{ex}

\begin{ex} %[2D11V55]
Một cái hồ hình chữ nhật có chiều dài $300\text{ m}$, chiều rộng $200\text{ m}$. Vận động viên chạy một vòng quanh bờ hồ hết bao nhiêu mét?
\choice
{$500\text{ m}$}
{$600\text{ m}$}
{$800\text{ m}$}
{\True $1 000\text{ m}$}
\loigiai{
Chu vi bờ hồ: $(300 + 200) \times 2 = 1 000$ m.
\textbf{Đáp án của câu hỏi là D}
}
\end{ex}

\begin{ex} %[2D11V55]
Trong các đơn vị đo độ dài mm, cm, dm, m, km, đơn vị đo nào lớn nhất?
\choice
{\True km}
{m}
{cm}
{dm}
\loigiai{
Ki-lô-mét (km) là đơn vị đo độ dài lớn nhất trong các đơn vị đã học.
\textbf{Đáp án của câu hỏi là A}
}
\end{ex}

\begin{ex} %[2D11V55]
Đoạn dây thứ nhất dài $45\text{ dm}$, đoạn dây thứ hai ngắn hơn đoạn thứ nhất $15\text{ dm}$. Cả hai đoạn dây dài bao nhiêu đề-xi-mét?
\choice
{$60\text{ dm}$}
{$70\text{ dm}$}
{\True $75\text{ dm}$}
{$80\text{ dm}$}
\loigiai{
Đoạn 2: $45 - 15 = 30$ dm. Cả hai đoạn: $45 + 30 = 75$ dm.
\textbf{Đáp án của câu hỏi là C}
}
\end{ex}

\textbf{Phần 2. Trắc nghiệm đúng sai}
\setcounterpageref{ex}{0}

\begin{ex} %[2D11V55]
Xét tính đúng sai của các phép đổi và quan hệ đơn vị đo độ dài:
\choiceTF
{\True 1 m = 100 cm.}
{\True 500 m + 500 m = 1 km.}
{1 km ngắn hơn 900 m.}
{\True Chiều cao của một học sinh lớp 2 khoảng 12 dm.}
\loigiai{
\begin{itemize}
    \item Ý a) Đúng: 1 m = 10 dm = 100 cm. Đúng.
    \item Ý b) Đúng: 1 000 m = 1 km. Đúng.
    \item Ý c) Sai: 1 km = 1 000 m > 900 m.
    \item Ý d) Đúng: 12 dm = 120 cm, phù hợp thực tế. Đúng.
\end{itemize}
\begin{itemize}
    \item a) \textbf{Đúng}.
    \item b) \textbf{Đúng}.
    \item c) \textbf{Sai}.
    \item d) \textbf{Đúng}.
\end{itemize}
}
\end{ex}

\begin{ex} %[2D11V55]
Cho ba đoạn thẳng: AB = 2 m, CD = 18 dm, EF = 250 cm:
\choiceTF
{\True AB dài hơn CD.}
{\True EF là đoạn thẳng dài nhất trong ba đoạn.}
{Tổng độ dài của AB và CD bé hơn 3 m.}
{\True Độ dài đoạn EF gấp 5 lần một đoạn dài 5 dm.}
\loigiai{
\begin{itemize}
    \item Ý a) Đúng: 2 m = 20 dm > 18 dm. Đúng.
    \item Ý b) Đúng: EF = 250 cm = 25 dm > 20 dm > 18 dm. Đúng.
    \item Ý c) Sai: 20 dm + 18 dm = 38 dm = 3 m 8 dm > 3 m.
    \item Ý d) Đúng: 25 dm : 5 = 5 dm. Đúng.
\end{itemize}
\begin{itemize}
    \item a) \textbf{Đúng}.
    \item b) \textbf{Đúng}.
    \item c) \textbf{Sai}.
    \item d) \textbf{Đúng}.
\end{itemize}
}
\end{ex}

\textbf{Phần 3. Trả lời ngắn}
\setcounterpageref{ex}{0}

\begin{ex} %[2D11V55]
Một chú ốc sên ban ngày bò lên được 8 dm, ban đêm tụt xuống 3 dm. Sau một ngày một đêm ốc sên bò lên được bao nhiêu đề-xi-mét?
\shortans{5}
\loigiai{
\textbf{Đáp số:} $5$\\
$8 - 3 = 5$ dm.
}
\end{ex}

\begin{ex} %[2D11V55]
Cuộn dây thép dài 1 km được chia thành các cuộn nhỏ dài 200 m. Được bao nhiêu cuộn nhỏ?
\shortans{5}
\loigiai{
\textbf{Đáp số:} $5$\\
$1\text{ km} = 1 000\text{ m}$. Số cuộn: $1 000 : 200 = 5$ cuộn.
}
\end{ex}

\begin{ex} %[2D11V55]
Thanh sắt dài 3 m được cưa thành các đoạn dài 5 dm. Cần cưa bao nhiêu lần để đứt hết?
\shortans{5}
\loigiai{
\textbf{Đáp số:} $5$\\
$3\text{ m} = 30\text{ dm}$. Số đoạn: $30 : 5 = 6$ đoạn. Cần $6 - 1 = 5$ lần cưa.
}
\end{ex}

\begin{ex} %[2D11V55]
Một hình tam giác có ba cạnh dài lần lượt là 20 cm, 30 cm, 50 cm. Chu vi của tam giác là bao nhiêu mét?
\shortans{1}
\loigiai{
\textbf{Đáp số:} $1$\\
$20 + 30 + 50 = 100\text{ cm} = 1$ m.
}
\end{ex}

\begin{ex} %[2D11V55]
Đoạn đường từ tỉnh A đến tỉnh B dài 45 km, từ tỉnh B đến tỉnh C dài 35 km. Đi từ A qua B đến C dài bao nhiêu ki-lô-mét?
\shortans{80}
\loigiai{
\textbf{Đáp số:} $80$\\
$45 + 35 = 80$ km.
}
\end{ex}

\begin{ex} %[2D11V55]
Một con mương dài 100 m. Người ta đã nạo vét được 60 m. Còn lại bao nhiêu mét mương chưa nạo vét?
\shortans{40}
\loigiai{
\textbf{Đáp số:} $40$\\
$100 - 60 = 40$ m.
}
\end{ex}
